\begin{frame}{El Desafío de las Altas Temperaturas}
  \frametitle{Requerimientos Térmicos en la Nueva Siderurgia}
  El desarrollo industrial impuso condiciones operativas que excedían la capacidad de los materiales convencionales.

  \begin{itemize}
    \item La producción de \textbf{acero} a escala industrial (ej. Proceso Bessemer) y la optimización de \textbf{máquinas de vapor} requerían hornos y calderas que operaran a temperaturas sostenidas y elevadas.
    \vspace{1cm}
    \item Los materiales de construcción existentes, como ladrillos comunes y mampostería, sufrían fallas estructurales y fusión bajo estas condiciones.
  \end{itemize}
\end{frame}

\begin{frame}{Aplicación de Cerámicos Refractarios}
  \frametitle{Soluciones para la Contención Térmica}
  La solución técnica se basó en el desarrollo de cerámicos especializados.

  \begin{itemize}
    \item Se industrializó la producción de ladrillos a partir de \textbf{arcillas refractarias (fireclay)} y \textbf{sílice ($SiO_2$)}, materiales con puntos de fusión muy superiores.
    \vspace{1cm}
    \item Su función principal fue el revestimiento interno de altos hornos, convertidores y calderas.
    \vspace{1cm}
    \item \textbf{Impacto técnico:} La capacidad de contener energía térmica de forma eficiente fue un requisito indispensable para la viabilidad de la siderurgia moderna.
  \end{itemize}
\end{frame}

\begin{frame}{Resistencia Química en Procesos Industriales}
  \frametitle{Contención de Agentes Corrosivos}
  La emergente industria química requería recipientes y ductos con alta inercia química.

  \begin{itemize}
    \item El \textbf{gres cerámico vitrificado} se estableció como el material estándar por su baja porosidad y alta resistencia a la corrosión por ácidos.
    \vspace{1cm}
    \item Se empleó en la fabricación de contenedores, tuberías y componentes para la producción de ácido sulfúrico y otros reactivos industriales.
  \end{itemize}
\end{frame}

\begin{frame}{Rol en la Infraestructura Sanitaria}
  \frametitle{Materiales para el Saneamiento Urbano}
  La rápida urbanización demandó soluciones duraderas para la gestión de aguas residuales.

  \begin{itemize}
    \item La crisis sanitaria en los centros industriales impulsó el desarrollo de redes de saneamiento a gran escala.
    \vspace{1cm}
    \item Las \textbf{tuberías de gres vitrificado} ofrecían la durabilidad, impermeabilidad y resistencia química necesarias para la construcción de los primeros sistemas de alcantarillado modernos.
  \end{itemize}
\end{frame}

\begin{frame}{El Problema del Aislamiento Eléctrico}
  \frametitle{Requerimientos para la Transmisión de Energía}
  La distribución de energía eléctrica presentaba un desafío fundamental: su aislamiento seguro y eficiente.

  \begin{itemize}
    \item Se necesitaba un material de \textbf{alta resistividad dieléctrica} para prevenir las pérdidas de corriente y los cortocircuitos.
    \vspace{1cm}
    \item Dicho material debía poseer \textbf{estabilidad mecánica y ambiental} para su uso en exteriores, soportando cargas mecánicas y exposición a la intemperie.
  \end{itemize}
\end{frame}

\begin{frame}{La Porcelana como Aislante Eléctrico}
  \frametitle{Material Dieléctrico Estándar}
  La cerámica, una vez más, proveyó la solución técnica óptima.

  \begin{itemize}
    \item La \textbf{porcelana eléctrica}, una cerámica densa a base de caolín, feldespato y cuarzo, demostró tener propiedades dieléctricas y mecánicas superiores.
    \vspace{1cm}
    \item Su principal aplicación fue la fabricación en masa de \textbf{aisladores} para las redes de telegrafía, telefonía y, fundamentalmente, para la transmisión de energía eléctrica.
  \end{itemize}
\end{frame}

\begin{frame}{Nuevos Requerimientos para la Electrónica}
  \frametitle{El Advenimiento de los Circuitos Eléctricos}
  El desarrollo de la radio y la telegrafía sin hilos introdujo nuevos requerimientos a nivel de componente.

  \begin{itemize}
    \item Los componentes electrónicos requerían materiales que ofrecieran \textbf{aislamiento eléctrico a pequeña escala}.
    \vspace{1cm}
    \item Adicionalmente, era necesaria una alta \textbf{estabilidad dimensional y térmica} para garantizar el funcionamiento fiable de los dispositivos.
  \end{itemize}
\end{frame}

\begin{frame}{Aislamiento en Tubos de Vacío}
  \frametitle{Componentes Críticos para Válvulas Termoiónicas}
  Los \textbf{tubos de vacío} fueron el componente activo central de la electrónica en la primera mitad del siglo XX.

  \begin{itemize}
    \item Las bases de estos dispositivos, que sujetaban los pines de conexión, debían ser excelentes aislantes y soportar el calor de operación.
    \vspace{1cm}
    \item Se fabricaron con \textbf{porcelana} y \textbf{esteatita}, asegurando el aislamiento entre electrodos y la integridad estructural del tubo.
  \end{itemize}
\end{frame}

\begin{frame}{Sustratos para Componentes Pasivos}
  \frametitle{Soportes Dieléctricos para Circuitos}
  La construcción de circuitos requería componentes pasivos como resistencias y capacitores.

  \begin{itemize}
    \item Durante este período, se utilizaron cerámicos como el \textbf{cuerpo o sustrato aislante} de dichos componentes.
    \vspace{1cm}
    \item El material cerámico proveía el soporte mecánico, el aislamiento eléctrico y la disipación de calor necesarios para su funcionamiento.
  \end{itemize}
\end{frame}

\begin{frame}{Puntos importantes}
  \begin{itemize}
    \item En el período analizado (1750-1950), la cerámica se consolidó como un \textbf{material de ingeniería} indispensable.
    \vspace{1cm}
    \item Si bien no fue una tecnología principal, funcionó como una \textbf{tecnología habilitadora fundamental}, operando como un componente sistémico clave.
    \vspace{1cm}
    \item Sus propiedades refractarias, de inercia química y de aislamiento dieléctrico fueron un requisito previo para el desarrollo de la siderurgia, la industria química y la electrificación a gran escala.
  \end{itemize}
\end{frame}
